\tituloI{METODOLOGÍA}

\tituloII{Método, tipo y alcance de la investigación}

\tituloIII{Método de investigación}
\parrafoIII{El método de investigación empleado es el **hipotético-deductivo**. Este método permite partir de una premisa teórica basada en la ingeniería de software y la norma ISO/IEC 25010, para luego someterla a comprobación mediante la observación empírica de métricas de código. A través de la deducción, se establecen las hipótesis que serán validadas o refutadas mediante el análisis estático de los resultados obtenidos tras la aplicación del Framework de Ingeniería de Prompts propuesto.}

\tituloIII{Tipo de investigación}
\parrafoIII{La investigación es de tipo **aplicada y tecnológica**. Es aplicada porque busca resolver un problema concreto de la industria del software: la acumulación de deuda técnica en sistemas generados por IA. Es tecnológica porque el aporte principal consiste en la creación de un Framework (herramienta metodológica) y un pipeline de evaluación automatizado que optimiza el proceso de producción de software, transformando el conocimiento teórico en una solución técnica ejecutable.}

\tituloIII{Alcance de la investigación}
\parrafoIII{De acuerdo con la Guía Metodológica, el alcance de esta investigación es **explicativo**. El estudio no se limita a describir cómo la IA genera código o a correlacionar variables, sino que pretende establecer una relación de causalidad. Se busca determinar cómo y por qué la estructuración sistémica de los prompts influye directamente en la reducción de la complejidad ciclomática y en la mejora de los índices de mantenibilidad del producto final.}

\tituloIII{Diseño de la investigación}
\parrafoIII{El diseño de la investigación es **experimental**, específicamente bajo un esquema **cuasi-experimental con pre-prueba y post-prueba**. Se utilizará un grupo de control donde la generación de código se realizará mediante métodos empíricos (prompts abiertos) y un grupo experimental donde se aplicará el Framework de Ingeniería de Prompts diseñado. Las mediciones de pre-prueba y post-prueba se realizarán utilizando herramientas de análisis estático (SonarQube/Radon) para comparar la evolución de la calidad estructural del software en ambos escenarios.}

\tituloII{Materiales y métodos}

\tituloIII{Materiales}
\parrafoIII{Los materiales empleados en la investigación se categorizan en recursos de hardware y software especializados para el desarrollo y análisis de sistemas:}

\begin{listaIII}
    \item \textbf{Recursos de Hardware:} Se utilizará una estación de trabajo con arquitectura x86\_64, procesador de múltiples núcleos y un mínimo de 16GB de RAM para garantizar la ejecución fluida de contenedores de análisis estático y modelos de lenguaje locales (clúster de pruebas bajo entorno Arch Linux).
    \item \textbf{Modelos de Lenguaje (LLMs):} Se emplearán modelos de acceso abierto como Llama 3 (8B/70B) y Mistral, seleccionados por su reproducibilidad y capacidad de procesamiento de lógica de programación.
    \item \textbf{Herramientas de Análisis Estático:} Se integrará \textit{SonarQube} (edición Community) para la medición de la deuda técnica e índices de mantenibilidad, y la librería \textit{Radon} para el cálculo automatizado de la Complejidad Ciclomática de McCabe.
    \item \textbf{Entorno de Desarrollo y Ejecución:} Uso de \textit{Docker} para la orquestación de servicios de métricas y \textit{Python/TypeScript} como lenguajes objetivo para la generación de código de microservicios.
\end{listaIII}

\tituloIII{Métodos}
\parrafoIII{El procedimiento se divide en fases secuenciales que aseguran el rigor del diseño cuasi-experimental:}

\begin{listaIIInum}
    \item \textbf{Fase de Línea Base (Pre-prueba):} Se ejecutará una serie de 50 peticiones de generación de código (prompts empíricos) sin estructura definida. Los resultados se procesarán a través del pipeline de análisis estático para establecer las métricas iniciales de mantenibilidad y deuda técnica.
    \item \textbf{Aplicación del Framework (Tratamiento):} Se aplicará el Framework de Ingeniería de Prompts diseñado, utilizando patrones arquitectónicos (como \textit{Chain-of-Thought} y estructuración basada en ISO 25010) para las mismas solicitudes de desarrollo de software.
    \item \textbf{Fase de Medición (Post-prueba):} El código resultante del grupo experimental será sometido nuevamente al pipeline automatizado. Se recolectarán los datos de complejidad, densidad de defectos y tiempo estimado de remediación.
    \item \textbf{Análisis Estadístico de Datos:} Se realizará una comparación de medias utilizando pruebas de significancia (como \textit{t-Student} o \textit{U de Mann-Whitney}) para validar las hipótesis planteadas, determinando si la mejora en la calidad estructural es estadísticamente significativa frente al método empírico.
\end{listaIIInum}